%! Characteristics of Exponential Functions — Unit 7, Lesson 7.0 — Exit Ticket
\documentclass[10pt]{article}
\usepackage{algebra2-article}
\usepackage{algebra2-boxes}
\newcommand{\TallMath}[1]{$\displaystyle #1\rule[-1.4em]{0pt}{3.2em}$}
\newcommand{\lgrid}[4]{%
  \draw[step=1, linegray!40, very thin] (#1,#3) grid (#2,#4);
  \draw[->, semithick, charcoal] (#1,0)--(#2,0) node[below right, font=\scriptsize]{$x$};
  \draw[->, semithick, charcoal] (0,#3)--(0,#4) node[left, font=\scriptsize]{$y$};}
\newcommand{\expcurve}[4]{\draw[very thick, forest, domain=#3:#4, samples=120, smooth]
  plot (\x,{#1*exp((#2)*(\x))});}

\begin{document}

\pageheader{Unit 7, Lesson 7.0}{Exit Ticket}
\namedateperiod

\vspace{0.08in}

No notes. The graph below is $f(x) = 8\left(\frac12\right)^{x}$. Use it for questions 1 and 2.
\begin{center}
\begin{tikzpicture}[scale=0.42]
  \lgrid{-1}{7}{-1}{10}
  \begin{scope}
    \clip (-1,-1) rectangle (7,10);
    \expcurve{8}{-0.693147}{-0.32}{7}
  \end{scope}
  \draw[navy, dashed, semithick] (-1,0)--(7,0);
  \node[navy, font=\scriptsize, above right] at (4.6,0) {$y=0$};
  \fill[charcoal] (0,8) circle (3pt) node[right, font=\scriptsize]{$(0,8)$};
  \fill[charcoal] (1,4) circle (3pt) node[right, font=\scriptsize]{$(1,4)$};
  \fill[charcoal] (2,2) circle (3pt) node[above right, font=\scriptsize]{$(2,2)$};
  \fill[charcoal] (3,1) circle (3pt) node[above right, font=\scriptsize]{$(3,1)$};
\end{tikzpicture}
\end{center}

\begin{enumerate}[leftmargin=1.8em, itemsep=8pt]
  \item \textbf{Read the features.} Write the domain and range in interval notation.

        $a=$ \blank{1.2cm} \quad $b=$ \blank{1.2cm} \quad Growth or decay? \blank{2.2cm}

        \vspace{3pt}
        Domain: \blank{3.0cm} \quad Range: \blank{3.0cm} \quad
        Horizontal asymptote: \blank{2.0cm}

        \vspace{3pt}
        $y$-intercept: \blank{1.8cm} \quad $x$-intercept: \blank{1.8cm} \quad
        $f$ is \blank{2.4cm} on its whole domain.

        \vspace{3pt}
        End behavior: as $x\to+\infty$, $f(x)\to$ \blank{1.4cm}; \; as $x\to-\infty$,
        $f(x)\to$ \blank{1.4cm}.

  \item \textbf{Explain.} Why does $f$ have no $x$-intercept? And why does the graph have no
        absolute minimum, even though it clearly flattens out along the bottom?
        \writelines{3}

  \item \textbf{SOL-style multiple choice.} What is the range of $g(x) = 6\left(\frac13\right)^{x}$?
        \begin{enumerate}[label=\Alph*., leftmargin=1.9em, itemsep=1pt]
          \item $(0,\infty)$
          \item $[0,\infty)$
          \item $(-\infty,\infty)$
          \item $(-\infty,0)$
        \end{enumerate}
\end{enumerate}

\end{document}
